% Generated from locale/tr/content/incompleteness/incompleteness-provability/provability-conditions.tex; do not hand-edit.


\olfileid[tr]{inc}{inp}{prc}

\olsection{$\Th{PA}$ için \usetoken{S}{derivability} Koşulları}

Peano aritmetiği, yani $\Th{PA}$, $\Th{Q}$'yu bütün !!{formula}s için
tümevarım aksiyomlarıyla genişleten kuramdır. Başka deyişle her
!!{formula}~$!A$ için $\Th{Q}$'ya
\[
(!A(0) \land \lforall[x][(!A(x) \lif !A(x'))]) \lif \lforall[x][!A(x)]
\]
biçimindeki aksiyomlar eklenir. Bunun aslında bir \emph{şema}, yani sonsuz
sayıda aksiyom olduğuna dikkat edin (nitekim $\Th{PA}$ {\em sonlu}
aksiyomlaştırılabilir değildir). Fakat bir simgeler dizisinin tümevarım
aksiyomunun bir örneği olup olmadığı etkin biçimde belirlenebildiğinden,
$\Th{PA}$'nın aksiyomlar kümesi hesaplanabilirdir. $\Th{PA}$, $\Th{Q}$'dan
çok daha güçlü bir kuramdır. Örneğin tümevarımı olağan biçimde kullanarak
toplamanın ve çarpmanın değişmeli olduğu kolayca ispatlanabilir. Hatta sonlu
yöntemlere dayanan sayı kuramsal ve kombinatorik argümanların çoğu $\Th{PA}$
içinde yürütülebilir.

$\Th{PA}$ hesaplanabilir biçimde aksiyomlaştırılmış olduğundan
$\Prf[\Th{PA}](x,y)$ !!{derivability} yüklemi hesaplanabilirdir ve dolayısıyla
$\Th{Q}$ içinde (öyleyse $\Th{PA}$ içinde de) temsil edilir. Daha önce olduğu
gibi bağıntıyı temsil eden formülü $\OPrf[\Th{PA}](x,y)$ ile göstereceğiz.
% [tr source emendation] Upstream uses the metalinguistic relation \Prf in
% this object-language definition; the representing formula is \OPrf.
$\OProv[\Th{PA}](y)$, sezgisel olarak ``$y$, $\Th{PA}$'nın aksiyomlarından
!!{derivable}{}dir'' diyen $\lexists[x][\OPrf[\Th{PA}](x,y)]$ formülü olsun.
$\Th{Q}$'nun aksiyomlarından biraz daha fazlasına gereksinmemizin nedeni,
kullandığımız kuramın birkaç temel olguyu !!{derive} edecek kadar güçlü
olduğunu bilmek zorunda olmamızdır. Bu olgular ilgili !!{derivability}
yüklemi hakkındadır.
Gereksindiğimiz olgular şunlardır:
\begin{enumerate}
\item[P1.] $\Th{PA} \Proves !A$ ise $\Th{PA} \Proves
  \OProv[\Th{PA}](\gn{!A})$ olur.
\item[P2.] Bütün $!A$ ve $!B$ !!{formula}s için
  \[
  \Th{PA} \Proves \OProv[\Th{PA}](\gn{!A \lif !B}) \lif
  (\OProv[\Th{PA}](\gn{!A}) \lif \OProv[\Th{PA}](\gn{!B})).
  \]
\item[P3.] Her $!A$ !!{formula} için
  \[
  \Th{PA} \Proves \OProv[\Th{PA}](\gn{!A})
  \lif \OProv[\Th{PA}](\gn{\OProv[\Th{PA}](\gn{!A})}).
  \]
\end{enumerate}
Bu üç özelliğin geçerli olduğunu doğrulamanın tek yolu
$\OProv[\Th{PA}](y)$ !!{formula}'nı dikkatle betimlemek ve ilgili biçimsel
!!{derivation}s{}i betimlemek için $\Th{PA}$'nın aksiyomlarını kullanmaktır.
(1) ve (2) koşulları kolaydır; asıl emek isteyen (3) koşuludur. (Bunun nasıl
bir emek gerektirdiğini düşünün \dots) Ayrıntıları tamamlamak zahmetli ve
ilgisiz olacağından burada $\Th{PA}$'nın yukarıda sıralanan üç özelliğe sahip
olduğunu kabul etmenizi isteyeceğiz. Makul bir $\OProv[\Th{PA}](y)$ seçimi
şunu da sağlar:
\begin{enumerate}
\item[P4.] $\Th{PA} \Proves \OProv[\Th{PA}](\gn{!A})$ ise
  $\Th{PA} \Proves !A$ olur.
\end{enumerate}
Fakat bu olguya gereksinmeyeceğiz.

\begin{digress}
Bu arada G\"odel de burada bizim olduğumuz kadar tembeldi. 1931 tarihli
makalesinin sonunda ikinci eksiklik teoreminin ispatını ana hatlarıyla verir
ve ayrıntıları sonraki bir makalede sunacağına söz verir. Buna hiçbir zaman
fırsat bulamadı; argümanı anlayan herkes ispatın tamamlanabileceğine inandığı
için ayrıntıları doldurması gerekmedi.
\end{digress}

